\scanpage{101}
% archive foliation: 103
2) $\uline{E'_a}$ Donc $\pi$ engendré par $u,v$ avec $\sigma u=u$, $\sigma v=u-v$
\[\sigma(x,y)=\sigma(xu+yv)=xu+y(u-v)=(x+y)u-yv.\]
\[\sigma(x,y)-(x,y)=yu-2yv=y(u-2v),\quad\text{donc la condition est}\]
$y\in\mathbb Z$, donc on trouve \deleted{que} les points de $\mathbb R\unsure{u_+|u_0}\bmod\pi$,
donc dans ce cas $X(\mathbb R)$ est connexe.

\margin{9.2.}
\uline{NB} Pour bien faire il faudrait aussi classifier les
courbes elliptiques réelles \unsure{intrinsèques}, i.e. les
torseurs (sur $\mathbb R$) sous les \inserted{schémas en} groupes elliptiques
$E_a,E'_a$. Les torseurs non triviaux fournissent les
surfaces conformes \deleted{sans} bord non orientables de
genre $1$. La suite exacte de $\gamma$-modules ($\gamma=\mathrm{Gal}(\mathbb C/\mathbb R)\simeq\{\pm1\}$)
\diagram{\begin{tikzcd}[ampersand replacement=\&]
0\arrow[r]\&\pi\arrow[r]\&V\arrow[r]\&E(\mathbb C)\arrow[r]\&0
\end{tikzcd}}
fournit
\[
H^1(\gamma,E(\mathbb C))\simeq H^2(\gamma,\pi)\simeq
\pi^\gamma/(1+\sigma)\pi=\pi_+/(1+\sigma)\pi
\]
1) Cas de $E_a$. On a $\pi\simeq\pi_+\oplus\pi_-$, d'où $\pi^\gamma=\pi_+$
\[
H^1(\mathbb R,E_{\mathbb R})\simeq H^1(\gamma,E(\mathbb C))\simeq H^2(\gamma,\pi)
\simeq\underset{\substack{\Vert\\\mathbb Z/2\mathbb Z}}{\uline{H^2(\gamma,\pi_+)}}
\oplus\underset{\substack{\Vert\\0}}{\uline{H^2(\gamma,\pi_-)}}
\]
Il y a exactement une classe non
nulle dans $H^1$, qui est d'ordre $2$.

2) Cas de $E'_a$. On a $H^1(\mathbb R,E_{\mathbb R})\simeq0$.

Donc les seules courbes elliptiques vraiment
\unsure{intrinsèques} sur $\mathbb R$ sont les $E_a^!=$ torseurs non
triviaux (d'ordre $2$) sous $E_a$, dépendant du
paramètre $a\in\mathbb R^{*+}$.
